\begingroup\small\setlength{\tabcolsep}{4pt}\renewcommand{\arraystretch}{1.18}
\begin{longtable}{>{\raggedright\arraybackslash}p{.21\linewidth}>{\raggedright\arraybackslash}p{.13\linewidth}>{\raggedright\arraybackslash}p{.59\linewidth}}
\caption{Avance documentado de los objetivos específicos; los nombres se abrevian y su formulación se conserva en el capítulo 2.}\label{editorial:avance}\\
\toprule\textbf{Objetivo} & \textbf{Estado} & \textbf{Evidencia y condición pendiente}\\\midrule\endfirsthead
\toprule\textbf{Objetivo} & \textbf{Estado} & \textbf{Evidencia y condición pendiente}\\\midrule\endhead
\bottomrule\endfoot
1. Identificar y caracterizar variables & Parcial & Capítulo~\ref{cap:caracterizacion}: auditoría y exploración ejecutadas. La selección base de A1.3 se explicita en la Sección~\ref{entrega:seleccionbase}; no depende de entrenar modelos. Quedan por acreditar la procedencia de la descarga y la ratificación del alcance académico.\\
2. Derivar indicadores & Parcial & Capítulo~\ref{cap:ingenieria} y Anexo~\ref{corr:anexodiccionario}: 26 candidatos y conjuntos por corte. A2.5 se limita a entregas y oportunidades; los puntajes permanecen en auditoría. A2.7 cuenta con bases exportadas, pendientes de acordar como entrada al modelado.\\
3. Comparar modelos & Parcial & La etiqueta temporal de A3.1 está implementada, con diferencias documentadas respecto del anteproyecto. El entrenamiento y la evaluación no se han iniciado; el Capítulo~\ref{editorial:modelos} delimita su estado.\\
4. Interpretar y seleccionar el modelo & No iniciado & Capítulo~\ref{editorial:shapestado}: SHAP y selección comparativa pendientes de modelos evaluados. El marco teórico no constituye un resultado de interpretación.\\
5. Diseñar el prototipo & No iniciado & Capítulo~\ref{editorial:prototipo}: pendiente de implementación. Las figuras explicativas de la tesis no equivalen al prototipo de predicciones e interpretabilidad previsto.\\
\end{longtable}\endgroup
Las diferencias de A2.5 y A3.1 requieren la ratificación académica ya señalada. La utilidad educativa se mantiene como propósito de comunicación y apoyo; no se afirma una mejora del aprendizaje o una reducción del abandono mediante intervenciones, cuya evaluación está excluida del alcance.
